\documentclass[12pt]{article}
\usepackage[utf8]{inputenc}
\usepackage[T1]{fontenc}
\usepackage{hyperref}
\title{Qatabanian Language Resources for K-12 Curriculum Development}
\author{Compiled by Claude Sonnet 5.5 (claude-sonnet-5-5) for Andr\'as Kornai}
\date{}
\begin{document}
\maketitle

\section{Introduction}
Qatabanian (ISO 639-3 xqt) is an extinct Ancient South Arabian language of the Afro-Asiatic family. It was used in the kingdom of Qatab\=an in what is now Yemen, mainly in Wadi Bayhan and Wadi Harib, and is attested from about 500 BC to 200 AD. About two thousand inscriptions survive in the Musnad script. Awsanian is closely related, but only about 25 inscriptions are known. The language has no speakers and, as far as I found, no school materials. Any use in schools would be through the inscriptions and the Musnad script, not through speaking.

\section{Primary Educational Resources}
I found no primers or textbooks for Qatabanian.

\section{Linguistic and Reference Materials}
\begin{itemize}
\item Qatabanian language (Wikipedia): \url{https://en.wikipedia.org/wiki/Qatabanian_language}
\item ``From inscriptions to grammar: notes on the grammar of non-Sabaic languages'' (Archaeopress), which discusses Qatabanian noun declension and the writing of vowels: \url{https://archaeopresspublishing.com/ojs/index.php/PSAS/article/view/1399}
\item Avanzini, \textit{Corpus of South Arabian Inscriptions I--III} (2004), Qatabanic, Marginal Qatabanic and Awsanite texts with editions and translations; the search result was a University of Pisa repository record: \url{https://arpi.unipi.it/handle/11568/87080}
\end{itemize}

\section{Religious and Community Resources}
None found. The inscriptions include religious and civic texts, but I did not search them for classroom use.

\section{Digital and Audio Resources}
The DASI database (Digital Archive for the Study of pre-Islamic Arabian Inscriptions) hosts epigraph records; these four records came up in the search and I did not examine them in detail:
\begin{itemize}
\item \url{https://dasi.cnr.it/project/1/epigraphs/947}
\item \url{https://dasi.cnr.it/project/1/epigraphs/948}
\item \url{https://dasi.cnr.it/project/1/epigraphs/5543}
\item \url{https://dasi.cnr.it/project/1/epigraphs/4949}
\end{itemize}
No audio exists, as the language is known only from inscriptions.

\section{Curriculum Development Context}
This is a classical, epigraphic language. A realistic school unit would cover Ancient South Arabian history, the Musnad script and selected inscriptions, using the DASI records and the Avanzini corpus as source material. I found no full grammar or dictionary of Qatabanian in this search.
\end{document}